\section{Experimental Results}

To prove the validity of our response-time analysis, we ran some
experiments on the software application on a real development board
(\textit{Olimex STM32\-H405}). In particular, the goal was to validate the
scheduling feasibility of the system using the computed workloads from the RTA,
and then obtain the values that make the system fail to meet its deadlines
(\textit{sensitivity analysis}). The latter experiments were only performed on
the Rust system. 

\subsubsection{Holistic Analysis}

From the holistic analysis, we obtained the workloads reported in Table
[\ref{tab:exp-holistic-baseline}] that maximize the system utilization. 
\begin{table}
  \centering
  \resizebox{\linewidth}{!}{%
    \begin{tabular}{lcc}
      \toprule
      \textbf{Task} & \textbf{Workload (ops)} & \textbf{Scaling} \\
      \midrule
      \texttt{Regular\_Producer}       & 15120 & 756 $\times$ 20 \\
      \texttt{On\_Call\_Producer}      & 10008 & 278 $\times$ 20 $\times$ 1.8 \\
      \texttt{Activation\_Log\_Reader} & 6755  & 139 $\times$ 20 $\times$ 1.8 $\times$ 1.35 \\
      \bottomrule
    \end{tabular}
  }
  \caption{Holistic baseline workloads (board experiments)}
  \label{tab:exp-holistic-baseline}
\end{table}
With this computed workloads both systems managed to meet all the deadlines
after execution of 50 \textit{hyperperiods}, which was considered to be a
sufficient number of tests.  

Then, we started experimenting with varying workloads by focusing on one task at
a time keeping the others fixed, we obtained the results in Table
[\ref{tab:exp-holistic-sensitivity}]. Workloads values were incremented at steps
of 100. 

\begin{table}
  \centering
  \resizebox{\linewidth}{!}{%
    \begin{tabular}{lccc l}
      \toprule
      \textbf{Scenario} & \textbf{RP (ops)} & \textbf{OCP (ops)} & \textbf{ALR (ops)} & \textbf{Observed} \\
      \midrule
      Baseline          & 15120 & 10008 & 6755  & No deadline misses \\
      \midrule
      RP increase       & 15200 & 10008 & 6755  & ALR misses $\sim$1/4 (when released with OCP) \\
      RP increase       & 16000 & 10008 & 6755  & RP always misses; ALR misses as before \\
      \midrule
      OCP increase      & 15120 & 10100 & 6755  & ALR misses $\sim$1/4 (when released with OCP) \\
      OCP increase      & 15120 & 16900 & 6755  & OCP always misses; ALR misses as before \\
      \midrule
      ALR increase      & 15120 & 10008 & 6900  & ALR misses $\sim$1/4 \\
      ALR increase      & 15120 & 10008 & 16900 & ALR always misses; no impact on others \\
      \bottomrule
    \end{tabular}
  }
  \caption{Holistic sensitivity results (board experiments)}
  \label{tab:exp-holistic-sensitivity}
\end{table}

Results matched closely the static analysis performed, since even small
increments from the original baseline incurred in deadline misses. These results 
also highlight a key property of FPS: under overload conditions the tasks gets
affected in priority order (from lower to highest). 

\subsection{Offset-based Analysis}

Using offset-based analysis we obtained the workloads in Table
[\ref{tab:exp-offset-baseline}] that maximize the system utilization. 

\begin{table}
  \centering
  \resizebox{\linewidth}{!}{%
    \begin{tabular}{lcc}
      \toprule
      \textbf{Task} & \textbf{Workload (ops)} & \textbf{Notes} \\
      \midrule
      \texttt{Regular\_Producer}       & 15120 & Baseline from analysis \\
      \texttt{On\_Call\_Producer}      & 15973 & Baseline from analysis \\
      \texttt{Activation\_Log\_Reader} & 1     & Baseline from analysis \\
      \bottomrule
    \end{tabular}
  }
  \caption{Offset-based baseline workloads (board experiments)}
  \label{tab:exp-offset-baseline}
\end{table}

No deadline misses were observed after 10 \textit{hyperperiods}.

The sensitivity analysis produced the values in Table [\ref{tab:exp-offset-sensitivity}].

\begin{table}
  \centering
  \resizebox{\linewidth}{!}{%
    \begin{tabular}{lccc l}
      \toprule
      \textbf{Scenario} & \textbf{RP (ops)} & \textbf{OCP (ops)} & \textbf{ALR (ops)} & \textbf{Observed} \\
      \midrule
      Baseline          & 15120 & 1 & 15973     & No deadline misses \\
      \midrule
      RP increase       & 16000 & 1 & 15973     & RP and ALR always miss \\
      \midrule
      OCP increase      & 15120 & 900   & 15973     & ALR misses $\sim$1/4 (when released with OCP) \\
      OCP increase      & 15120 & 16900 & 15973     & OCP always misses; ALR misses as before \\
      \midrule
      ALR increase      & 15120 & 15973 & 16900 & ALR always misses; no impact on others \\
      \bottomrule
    \end{tabular}
  }
  \caption{Offset-based sensitivity results (board experiments)}
  \label{tab:exp-offset-sensitivity}
\end{table}

This time, results were slightly more distant from the ones obtained from static
analysis. This could be due to an overestimation of the WCET of system and
operation overhead. 

Nonetheless, results still follow a consistent behavior when overload is
reached, with the lowest priority jobs suffering from interference and incurring
in deadline misses. 

